\subsection{A share history changes the denominator}
\label{sec:eps-lifecycle}

A company can earn the same profit with different numbers of shares
outstanding during the year. Dividing by the closing count would attribute
the entire year's earnings to shares that may have existed for only a few
months. Basic earnings per share instead relates earnings attributable to
ordinary shareholders to the time those shares were outstanding. Diluted
earnings per share asks a further, conditional question: how would that
measure change if eligible potential shares had become ordinary shares?
Neither calculation exercises an option or converts a bond in the company's
actual accounts.

The following synthetic case uses IAS~33 in the retained original EU
Regulation 2023/1803, including its integral Appendix~A. It develops a method
under that historical text; applying it to a real issuer requires the
issuer's reporting basis, operative version and instrument terms. It is
separate from the Apple disclosures that follow
\citep[IAS~33, paragraphs 2--4, 9--12, 30--36]{eu20231803}.
Assume one fully paid ordinary-share class, one consolidated reporting
population, and a 365-day year from 1 January to 31 December 2025.
Monetary amounts are currency units; share counts are numbers of shares,
so EPS is currency units per share.

The year opens with 1,000 ordinary shares and no treasury holdings.
On 1 April, consideration becomes receivable for 200 new shares at 12
each and is paid that day: cash and contributed equity rise by 2,400.
On 1 July, a two-for-one split doubles every outstanding ordinary share.
On 1 October, the company buys 300 shares for 7 each and holds them in
treasury. Cash and equity fall by 2,100; issued shares remain 2,400,
treasury shares become 300, and outstanding shares become 2,100.
The split contributes no additional resources. IAS~33 therefore restates
earlier share counts onto the new share basis, whereas the resource-changing
issue and repurchase enter the average from their respective dates
\citep[IAS~33, paragraphs 19--21, 26--29]{eu20231803}.

\begin{center}
\begin{tabular}{lrrr}
\toprule
Interval, including first date & Days & Actual outstanding & Split-adjusted\\
\midrule
1 January--1 April & 90 & 1,000 & 2,000\\
1 April--1 July & 91 & 1,200 & 2,400\\
1 July--1 October & 92 & 2,400 & 2,400\\
1 October--1 January 2026 & 92 & 2,100 & 2,100\\
\bottomrule
\end{tabular}
\end{center}

Each interval excludes its end date, so no day is counted twice. The basic
denominator is consequently
\[
 D=\frac{90(2000)+91(2400)+92(2400)+92(2100)}{365}
   =\frac{162480}{73}\approx2225.7534.
\]
The closing count of 2,100 answers a different question. Nor should the
split-adjusted count overwrite the historical register: there really were
1,000 outstanding shares in January. If the comparative year had constant
shares of 1,000 and ordinary earnings of 6,000, its previously calculated
EPS of 6 would be presented as 3 after this split. The same retrospective
principle applies to a split occurring after year-end but before
authorization of the statements. A later cash issue does not receive that
treatment \citep[IAS~33, paragraphs 64--65, 70(d)--71]{eu20231803}.

\subsection{Attribute earnings before testing potential shares}

Suppose consolidated profit is 15,000 from continuing operations and
18,000 in total. Non-controlling interests receive 1,000, all from
continuing operations. A cumulative, nonparticipating preference class,
classified here as equity, has a current-year entitlement of 2,000.
Deduct that entitlement even though it has not been declared. A payment
of 500 of prior-year arrears is not another deduction from this year's
earnings. These are supplied attribution facts; the EPS calculation does
not infer them from cash dividends
\citep[IAS~33, paragraphs 12--14 and A1]{eu20231803}.

Earnings attributable to the parent's ordinary shareholders are thus
$15000-1000-2000=12000$ from continuing operations and
$18000-1000-2000=15000$ in total. Discontinued operations contribute
3,000. Basic continuing EPS is $12000/D\approx5.3914$; basic total EPS
is $15000/D\approx6.7393$. Their difference, approximately 1.3479,
is basic discontinued EPS. Profit already includes the interest and tax
effects described below; adding interest again at this stage would be wrong.

There are three potential-share series. First, cash options have been
outstanding all year: 200 shares at an exercise price of 4 before the
split, and 400 at 2 afterward. Use a supplied average market price of 5
on the post-split basis. Earlier prices must be converted to that same
basis. The options require no future employee services and contain no
performance condition, debt tender, redemption requirement or cash-settlement
alternative.

Assumed exercise would provide $400(2)=800$. At the average market price,
800 corresponds to 160 fairly priced shares. Only the remaining
$400-160=240$ shares enter the diluted denominator; there is no earnings
adjustment. Adding all 400 shares would ignore the assumed proceeds.
IAS~33 describes this as the difference between assumed issue and the
shares represented by proceeds at the average price
\citep[IAS~33, paragraphs 45--47 and A4--A5]{eu20231803}.
The calculation records neither cash of 800 nor an actual share issue.

Second, a convertible debt series is issued on 1 July. Its original terms
permit voluntary conversion into 600 post-split ordinary shares; conversion
is neither mandatory nor subject to a cash alternative. For this case,
supplied issue proceeds of 10,000 divide into a liability of 9,200 and
conversion equity of 800. The recognized effective-interest expense for
the remaining year is 800, with a related tax benefit of 200 and no other
earnings consequences. The original terms pay a coupon of 400 on 1 October
and principal of 9,200 plus a final coupon of 400 on 1 January 2026.
The supplied effective rate is $1/23$ per 92-day period: discounting those
payments gives $400/(24/23)+9600/(24/23)^2=9200$.
Interest of 400 first raises the liability to 9,600; the October coupon
returns it to 9,200, and another 400 of interest gives a year-end liability
of 9,600 before January settlement. The intervening payment explains why
the two interest amounts are equal. These supplied amounts include the
relevant cost and discount effects; substituting a coupon alone need not do so.
Instrument classification and component valuation are separate accounting
inputs, illustrated by the compound-bond example below.

Assumed conversion from issue eliminates the net expense of
$800-200=600$ and adds shares only for the 184 days that the instrument
exists. Its contributions are
\[
 \Delta N_b=600,\qquad
 \Delta D_b=600\frac{184}{365}=\frac{22080}{73},
 \qquad
 \frac{\Delta N_b}{\Delta D_b}=\frac{365}{184}\approx1.9837.
\]
IAS~33 requires the after-tax interest effect and any consequential
earnings changes, together with this issue-date weighting
\citep[IAS~33, paragraphs 33--39, 49--50]{eu20231803}.
For example, a profit-sharing expense caused by removing interest would
also change the numerator; it cannot disappear merely because a coupon
and a tax rate are available.

Third, the preference class can convert into 400 post-split ordinary
shares and has existed throughout the year. Its assumed conversion adds
back the current entitlement of 2,000. Earnings per incremental share
are therefore $2000/400=5$. The class has no participation in residual
earnings; participating rights would require a different attribution
calculation.

\subsection{The order of dilution changes the answer}

Continuing-operations earnings attributable to ordinary shareholders,
12,000 here, control admission. Each series is assessed separately in
ascending order of earnings per incremental share: options at zero,
debt at about 1.9837, then preference shares at 5. Recalculate after each
admitted series; comparing every instrument only with basic EPS loses
this dependency \citep[IAS~33, paragraphs 41--44]{eu20231803}.

The options reduce continuing EPS from approximately 5.3914 to 4.8667.
The debt then reduces it further:
\[
 D_{\mathrm{dil}}=D+240+\frac{22080}{73}
                 =\frac{202080}{73},
 \qquad
 E_{\mathrm{dil,cont}}
   =\frac{12000+600}{D_{\mathrm{dil}}}
   =\frac{7665}{1684}\approx4.5517.
\]
The preference class would increase this result because its incremental
ratio, 5, exceeds 4.5517. Exclude it. Yet $5<5.3914$, so it would have
appeared dilutive if tested against basic EPS alone. That is why the
order matters.

The admitted population also determines the other EPS amounts. Total
earnings increase from 15,000 to 15,600 through the same debt addback;
discontinued earnings remain 3,000. Thus
\[
 E_{\mathrm{dil,total}}=\frac{15600}{202080/73}
       =\frac{4745}{842}\approx5.6354,\qquad
 E_{\mathrm{dil,disc}}=\frac{3000}{202080/73}
       =\frac{1825}{1684}\approx1.0837.
\]
Continuing and discontinued EPS sum to total EPS before presentation
rounding. The report retains that reconciliation, each admitted numerator
and denominator increment, and the excluded preference series. Basic and
diluted EPS receive equal prominence; negative EPS is still presented
\citep[IAS~33, paragraphs 66--70]{eu20231803}.

Two sign changes test the control number. If continuing ordinary earnings
are instead a loss of 5,000 while total ordinary earnings remain a profit
of 15,000, none of these three series is admitted: each would reduce the
loss per share from continuing operations. Total profit cannot reverse
that decision. Conversely, retain continuing earnings of 12,000 but
change total ordinary earnings to a loss of 9,000. Options and debt still
enter, so diluted total EPS uses $-9000+600=-8400$ over the same diluted
denominator. It is less negative than basic total EPS. That is the
consequence of the continuing-operations control, illustrated explicitly
in the standard's Appendix~A
\citep[IAS~33, paragraphs 41--43 and A3]{eu20231803}.

\subsection{Actual conversion, expiry and incomplete terms}

Actual events divide the year differently. Suppose the debt pays its coupon
and then converts under its original terms on 1 October; interest ceases
that day. Before conversion, recognized interest is 400 with a tax benefit
of 100. Paying the coupon reduces the liability to 9,200, which transfers
to ordinary equity; the original conversion-equity component remains in equity.
After conversion, 600 actual shares enter basic EPS for 92 days.
The diluted calculation adds only the preceding 92 days of potential
shares and the preconversion net expense of 300. There are therefore
$600(92+92)/365$ additional share-years in the diluted comparison, with
no overlap. Under the otherwise identical supplied earnings assumptions,
basic ordinary earnings rise by the 300 of avoided later expense, and
the remaining addback of 300 restores the same fully diluted numerator
as before. Treating the potential shares as outstanding for all 184 days
\emph{in addition} to the actual shares would double count
\citep[IAS~33, paragraphs 21(c), 36--38]{eu20231803}.

Expiry has no subsequent actual-share leg. If the options instead lapse
on 1 October, their contribution is $240(273/365)$, using a supplied
average price of 5 for this variant's relevant period. They do not vanish
from the earlier part of the year. An exercise price equal to or above
the average price gives zero incremental shares in this simple option
profile. More generally, a conversion that leaves continuing EPS exactly
unchanged is not dilutive: paragraph 41 requires a decrease in profit
per share or an increase in loss per share.

Employee awards require another input even when their other terms fit
the option calculation. Under paragraph 47A, assumed proceeds include
the IFRS~2 value of future goods or services. If that supplied value were
400 across these 400 options, the assumed price would be
$2+400/400=3$, and the increment at market price 5 would be 160 shares.
This conditional arithmetic neither values the services nor calculates
award expense. Performance-based awards follow the contingent-share
provisions instead; share-or-cash alternatives, written puts, rights
issues and participating or partly paid shares also need their own
methods \citep[IAS~33, paragraphs 47A--63, A2 and A13--A16]{eu20231803}.

An EPS assessment therefore references a dated book, reporting population,
share class, attributed earnings, share history, prices and complete
instrument terms. It produces report amounts and their reconciliation,
without changing the book. Missing average prices, future-service values,
conversion consequences or settlement terms leave diluted EPS unresolved
even when basic EPS is available. Missing attribution or ordinary-share
timing can prevent both. Known complex terms select the relevant further
method; they are not grounds for silently dropping the instrument.
These distinctions let the following issuer disclosures supply facts
without being mistaken for a complete method or a complete history.
